\section{1 de Julio}
\label{sec:1-de-julio}

En estos párrafos voy comentar que hay hecho hasta ahora en el proyecto. El objetivo es construir un emulador que sea capaz de ejecutar algo de GBA. La idea inicial de hacer este trabajo se esboza más considerablemente a inicios del 2026, en plena época de exámenes finales como un punto de partida para responder a la pregunta de ¿y qué hago después de estos examenes? o atacar a los pensamientos del estilo: ¡Qué viene el TFG!. Años anteriores ya tenía una muy muy y mucho vaga idea de querer hacer un emulador. Me parece una pieza del software moderno muy única en sus objetivos. Algo que parece que vaya incluso en contra de lo que pida hoy en día. Quizás hubiera sido más fácil (ojo, o más complejo) haber optado por hacer una página web, entrenar una ia (o llm), crear unos materiales, un estudio sobre algún acontecimiento de interés. Sin embargo, desde antes considero que emulación es un terrno fangoso que se prefiere evitar. Siempre me ha parecido muy curioso poder ejecutar programas compatibles de una máquina en otra que sean incompatibles. Realmente, llegar a esa situación es complejo, porque implica conocer al detalle el funcionamiento de una máquina. Y más si tenemos en cuenta que la documentación de esa máquina a la que pretendemos emular sea opcaco o inexistente. Esto conlleva a tratar al sistema objetivo de emular como una caja negra y por ello a aplicar técnicas de ingiería inversa. La emulación es un aspecto que suele estar relacionado a la retrocomputación y en concreto a los videojuego porque promueven la preservación de obras eficazmente. La emulación también suele ser usada para pruebas y depueración de aplicaciones y productos. Es por este objeto que los fabricantes de hardware son los primeros interesados en proporcionar las facilidades a sus interesados en que su software funciona bien. Los kits de desarrollo porporcionan unidades de pruebas y/o emuladores/simuladores para dichas pruebas. Por ejemplo: kits de desarrollo, android studio... Otros usos a los que se suele dar la emulación o un concepto muy similar es: resetta 2 (programa de traducción de mac de intel a arm), wine (x86 windows en x86 linux).

Como es habitual, las videoconsolas más antiguas tienden a ser objeto de estudio para los cimientos informáticos. Una videoconsola retro tiene la ventaja de gozar con una comunidad apasionada por seguir redescubriendo las entrañas y los títulos para estas. Además de una amplia cultura popular retro arraigada en la sociedad que ha atraido una parte importante del público directamente o indirectamente a conocerla. En lo que atañe a los emuladores de consolas, hoy gozan de una aceptación para la preservaación del videojuego al cumplir años de mejoras y de consolidación que facilitan una experiencia estable y más recreable al ejecutar. Los emuladores suelen ser de código abierto (es demasiado raro ver que sean de pago con excepciones de las mismas empresas o compañías con poco reconociemineto que aplican prácticas dudosas). Por ello, mi motivación no es reinventar la rueda, ni tampoco mejorarla, sino plantearlo como una cuestión de aprendizaje. Pienso que es interesante conocer de primera mano los mecanismos de ejecución de programas antiguos. Así como sus optimizaciones y mejoras sobre los programas. Me he documentado y descargado manuales de arm, documentación no ofcial e investigado proyectos.

La GBA tiene dos chips pero solo haremos uno porque es el se usa para juegos de gba, el otro es solo para gb. La cpu encargada de ejecutar esto del gba es arm7tdmi. Este procesador es el típico ejemplo de estudio de la arquitectura arm para dar a conocer sus registros bancados, modos de operación e implementación de dos conjuntos de instrucciones diferentes (propiamente dichos son arm y thumb). El objetivo de las thumb es la mejora de rendimiento ofreciendo mayor throught en el procesamiento de instrucciones puesto que ocupan la mitad de una de arm.

Inicialmente, empecé haciendo un \enquote{prototipo} en c++ idiomático que \enquote{lograba} sumar dos números. Pero no me convencía del todo escribir todo tan idomático porque se hacía bola y reinicié el proyecto. La idea del proyecto es ir implementando por partes y con pruebas. Para esto utilizaremos el gtest como framework. A la hora de volver a realizar la implementación he decido empezar por la CPU ARM7TDMI, pues resulta en el núcleo fundamente que tiene esta videoconsola. La decisión de realizar este componente primero reside en un motivo lógico de cómo ir estructurando los diferentes componentes. (¿Qué va primero el sistema operativo o las aplicaciones para este?). Pese a que pueda ser lógico pudiese ser útil adoptar otro punto de vista en el empecemos por las aplicaciones en lugar del sistema operativo. Similiar a como sucedió la combinación GNU/Linux, donde GNU tenía hecho el conjunto de utilidades para ser usadas antes pero no el núcleo y en paralelo el kernel Linux ya estaba hecho.  

que contar:
- motivación y objetivos.
- cómo edito, cómo documento, repo git, documentación.
- marco teórico: (gba, contexto), c++ (testing, editor, framework, biblioteca)
- c-style c++ implica limitarse -> pero idomatico hard -> transicion poco.
- decidí implementar extractores y constructores de instrucciones según docs.
- primero con macros pero endeble. luego templates mejor -> tests iban más rápido.
- registros -> ¿cómo es la estructura? -> ¿porque cambiarla? -> macros
- flags de registros -> lógica booleana implica conocer
- planteamiento inicial de métodos de cpu.
- las pruebas son vitales para comprobar que funciona sin tener un programa completo. -> detección de missspells -> util para comporbar que hace bien el cambio de modo.

%%% Local Variables:
%%% mode: LaTeX
%%% TeX-master: "../main"
%%% End:
